% Part: first-order-logic
% Chapter: sequent-calculus
% Section: proof-theoretic-notions

\documentclass[../../../include/open-logic-section]{subfiles}

\begin{document}

\begin{editorial}
  Esta seção reúne as definições da relação de demonstrabilidade e de
  consistência para a dedução natural.
\end{editorial}

\iftag{FOL}
      {\olfileid[pt]{fol}{seq}{ptn}}
      {\olfileid[pt]{pl}{seq}{ptn}}

\olsection{Noções da Teoria da Prova}

\begin{explain}
Assim como definimos uma série de noções semânticas importantes
(validade, consequência, satisfatibilidade), definimos agora as
\emph{noções da teoria da prova} correspondentes.  Elas não são definidas
por apelo à satisfação de !!{sentence}s em !!{structure}s, mas por apelo à
!!{derivability} ou à !!{nonderivability} de certos sequentes.  Foi
uma descoberta importante que essas noções coincidem.  Que elas coincidem é
o conteúdo dos teoremas da \emph{correção} e da \emph{completude}.
\end{explain}

\begin{defn}[Teoremas]
!!^a{sentence}~$!A$ é um \emph{teorema} se há !!a{derivation}
em~$\Log{LK}$ do sequente $\quad \Sequent !A$.  Escrevemos $\Proves
!A$ se $!A$ é um teorema e $\Proves/ !A$ se não é.
\end{defn}

\begin{defn}[!!^{derivability}]
!!^a{sentence}~$!A$ é \emph{!!{derivable} de} um conjunto de
!!{sentence}s~$\Gamma$, $\Gamma \Proves !A$, se, e somente se, há um
subconjunto finito~$\Gamma_0 \subseteq \Gamma$ e uma sequência $\Gamma_0'$
das !!{sentence}s em~$\Gamma_0$ tal que $\Log{LK}$ !!{derive}
$\Gamma_0' \Sequent !A$.  Se $!A$ não é !!{derivable} de $\Gamma$,
escrevemos $\Gamma \Proves/ !A$.
\end{defn}

Por causa das regras de contração, enfraquecimento e troca, a ordem
e o número de !!{sentence}s em~$\Gamma_0'$ não importam: se um
sequente $\Gamma_0' \Sequent !A$ é !!{derivable}, então também o é
$\Gamma_0'' \Sequent !A$ para qualquer $\Gamma_0''$ que contenha as mesmas
!!{sentence}s que~$\Gamma_0'$.  Por exemplo, se $\Gamma_0 = \{!B, !C\}$,
então tanto $\Gamma_0' = \tuple{!B, !B, !C}$ quanto $\Gamma_0'' =
\tuple{!C, !C, !B}$ são sequências que contêm apenas as !!{sentence}s
em~$\Gamma_0$. Se um sequente contendo uma delas é !!{derivable}, o
outro também é, por exemplo:
\begin{prooftree}
  \AxiomC{}
  \Deduce$!B, !B, !C \fCenter !A$
  \RightLabel{\LeftR{\Contraction}}
  \UnaryInf$!B, !C \fCenter !A$
  \RightLabel{\LeftR{\Exchange}}
  \UnaryInf$!C, !B \fCenter !A$
  \RightLabel{\LeftR{\Weakening}}
  \UnaryInf$!C, !C, !B \fCenter !A$
\end{prooftree}
De agora em diante, diremos que, se $\Gamma_0$ é um conjunto finito de
!!{sentence}s, então $\Gamma_0 \Sequent !A$ é qualquer sequente em que o
antecedente é uma sequência de !!{sentence}s em~$\Gamma_0$, e incluiremos
tacitamente contrações, trocas e enfraquecimentos se necessário.

\begin{defn}[Consistência]
Um conjunto de sentenças~$\Gamma$ é \emph{inconsistente} se, e somente se, há um
subconjunto finito~$\Gamma_0 \subseteq \Gamma$ tal que $\Log{LK}$
!!{derive} $\Gamma_0 \Sequent \quad$. Se $\Gamma$ não é
inconsistente, isto é, se para todo $\Gamma_0 \subseteq \Gamma$ finito,
$\Log{LK}$ não !!{derive} $\Gamma_0 \Sequent \quad$, dizemos que é
\emph{consistente}.
\end{defn}

\begin{prop}[Reflexividade]
\ollabel{prop:reflexivity}
Se $!A \in \Gamma$, então $\Gamma \Proves !A$.
\end{prop}

\begin{proof}
O sequente inicial $!A \Sequent !A$ é !!{derivable}, e $\{!A\}
\subseteq \Gamma$.
\end{proof}
  
\begin{prop}[Monotonicidade]
\ollabel{prop:monotonicity}
Se $\Gamma \subseteq \Delta$ e $\Gamma \Proves !A$, então $\Delta
\Proves !A$.
\end{prop}

\begin{proof}
Suponha $\Gamma \Proves !A$, isto é, há um $\Gamma_0
\subseteq \Gamma$ finito tal que $\Gamma_0 \Sequent !A$ é
!!{derivable}. Como $\Gamma \subseteq \Delta$, então $\Gamma_0$ também
é um subconjunto finito de~$\Delta$. A !!{derivation} de $\Gamma_0
\Sequent !A$ mostra, assim, também que $\Delta \Proves !A$.
\end{proof}

\begin{prop}[Transitividade]
\ollabel{prop:transitivity}
Se $\Gamma \Proves !A$ e $\{!A\} \cup \Delta \Proves
!B$, então $\Gamma \cup \Delta \Proves !B$.
\end{prop}

\begin{proof}
Se $\Gamma \Proves !A$, há um $\Gamma_0 \subseteq \Gamma$ finito
e !!a{derivation}~$\pi_0$ de $\Gamma_0 \Sequent !A$.  Se $\{!A\}
\cup \Delta \Proves !B$, então, para algum subconjunto finito $\Delta_0
\subseteq \Delta$, há !!a{derivation}~$\pi_1$ de $!A, \Delta_0
\Sequent !B$. Considere a seguinte !!{derivation}:
\begin{prooftree}
  \AxiomC{}
  \RightLabel{$\pi_0$}
  \Deduce$\Gamma_0 \fCenter !A$
  \AxiomC{}
  \RightLabel{$\pi_1$}
  \Deduce$!A, \Delta_0 \fCenter !B$
  \RightLabel{\Cut}
  \BinaryInf$\Gamma_0, \Delta_0 \fCenter !B$
\end{prooftree}
Como $\Gamma_0 \cup \Delta_0 \subseteq \Gamma \cup \Delta$, isso
mostra $\Gamma \cup \Delta \Proves !B$.
\end{proof}

Note que isso significa, em particular, que se $\Gamma \Proves !A$ e $!A
\Proves !B$, então $\Gamma \Proves !B$. Segue-se também que se $!A_1,
\dots, !A_n \Proves !B$ e $\Gamma \Proves !A_i$ para cada~$i$, então
$\Gamma \Proves !B$.

\begin{prop}
\ollabel{prop:incons}
$\Gamma$ é inconsistente se, e somente se, $\Gamma \Proves {!A}$ para toda
  sentença~$!A$.
\end{prop}

\begin{proof}
Exercício.
\end{proof}

\begin{prob}
Prove \olref[fol][seq][ptn]{prop:incons}
\end{prob}

\begin{prop}[Compacidade]
  \ollabel{prop:proves-compact}
  \begin{enumerate}
  \item Se $\Gamma \Proves !A$, então há um subconjunto finito $\Gamma_0
    \subseteq \Gamma$ tal que $\Gamma_0 \Proves !A$.
  \item Se todo subconjunto finito de~$\Gamma$ é
    consistente, então $\Gamma$ é consistente.
  \end{enumerate}
\end{prop}

\begin{proof}
  \begin{enumerate}
    \item Se $\Gamma \Proves !A$, então há um subconjunto finito
      $\Gamma_0 \subseteq \Gamma$ tal que o sequente $\Gamma_0
      \Sequent !A$ tem !!a{derivation}. Consequentemente, $\Gamma_0
      \Proves !A$.
    \item Se $\Gamma$ é inconsistente, há um subconjunto
      finito~$\Gamma_0 \subseteq \Gamma$ tal que $\Log{LK}$
      !!{derive} $\Gamma_0 \Sequent \quad$. Mas então $\Gamma_0$ é um
      subconjunto finito de~$\Gamma$ que é inconsistente.
  \end{enumerate}
\end{proof}

\end{document}
